% !PRINTED 576
\weberpage{576}
%<<CONT
nach der Bedeutung (12) von $\Pi(t)$ folgt:
\[ \Pi'(\tau) \equiv (\tau-\tau_1)(\tau-\tau_2) \cdots (\tau-\tau_{f-1}) \equiv 0 \pmod{\mathfrak p}. \]
Dies ist aber unmöglich, weil die $\tau, \tau_1, \ldots \tau_{f-1}$, wie wir gesehen haben, nach dem Modul $\mathfrak p$ incongruent sind. Damit ist also unser Satz 4. vollständig bewiesen. Wir können hiernach, wenn $x, y$ zwei neue Variable bedeuten, ein Functional $\pi$ des Ideals $\mathfrak p$ bilden:
\weq{23}{\pi = x p + y P(\tau).}
Wir betrachten nun ganze Formen $\Phi(t)$ in $R$, die durch die Substitution $t=\tau$ nicht nur durch $\mathfrak p$, sondern durch die natürliche Primzahl $p$ theilbar werden, also der Congruenz
\weq{24}{\Phi(\tau) \equiv 0 \pmod p}
genügen. Die Primzahl $p$ möge folgendermaassen in ihre Primfactoren in $\Omega$ zerlegt sein:
\weq{25}{p = \mathfrak p \mathfrak p_1 \mathfrak p_2 \cdots,}
worin
\[ \mathfrak p, \mathfrak p_1, \mathfrak p_2, \ldots \]
gleiche oder verschiedene Primideale der Grade
\[ f, f_1, f_2, \ldots \]
sind. Diesen Primfactoren entspricht (nach dem Satze 3.) eine Reihe ganzer rationaler Formen
\[ P(t), P_1(t), P_2(t), \ldots \]
der Grade $f, f_1, f_2, \ldots$, und wenn etwa $\mathfrak p$ mit $\mathfrak p_1$ identisch ist, so ist auch $P(t)$ mit $P_1(t)$ identisch. Wenn man in (25) rechts und links die Norm nimmt, und die Formel $N(\mathfrak p) = p^f$ berücksichtigt [\S.\,140, (3), \S.\,151, (3)], so folgt
\weq{26}{n = f + f_1 + f_2 + \cdots}
Wenn nun $\Phi(t)$ eine der Bedingung (24) genügende ganze rationale Form ist, so folgt aus dem Satze 3.
\weq{27}{\Phi(t) \equiv P(t) \Phi_1(t) \pmod p,}
worin $\Phi_1(t)$ eine ganze Function in $R$ ist, die der Bedingung
\[ P(\tau) \Phi_1(\tau) \equiv 0 \pmod{\mathfrak p \mathfrak p_1 \mathfrak p_2 \ldots} \]
genügt. Nach dem Satze 4. folgt hieraus
\weq{28}{\Phi_1(\tau) \equiv 0 \pmod{\mathfrak p_1 \mathfrak p_2 \ldots},}
%CONT>>
